% Generated by Listas/make_statements.py from pratica-lecture-01.md, item 9. Do not edit.
\textbf{Source.} Tadelis (2013), \emph{Game Theory: An Introduction}, Exercise 3.6 ``Price Competition'', p.~58 {[}PDF 77{]}. Reproduced verbatim. Answer in the source: none in the book (the preface, p.~xiii, says about half are solved on Princeton UP's website; not checked).



\begin{quote}

\textbf{Price Competition:} Imagine a market with demand \(p(q) = 100 - q\). There are two firms, 1 and 2, and each firm \(i\) has to simultaneously choose its price \(p_i\). If \(p_i < p_j\), then firm \(i\) gets all of the market while no one demands the good of firm \(j\). If the prices are the same then both firms split the market demand equally. Imagine that there are no costs to produce any quantity of the good. (These are two large dairy farms, and the product is manure.) Write down the normal form of this game.

\end{quote}

